% registry: BP
\methodhead{BP}{B\"ohme and Papuga (2024): data-driven integral criterion}\label{sheet:BP}

\block{Source} S.~A.~B\"ohme, J.~Papuga, \emph{Fatigue Fract.\ Eng.\ Mater.\ Struct.} 47 (2024)
\cite{bohme2024}, Eq.~(16). The equations and coefficients implemented are those restated by
B\"ohme, Papuga and Lange (2026) \cite{bohme2026}, Eqs.~(1), (2), (6), (9), (17); Tables~1, 2, 4.

\block{Idea and scope} Stress-based integral criterion. Instead of selecting one critical plane, the
squared shear and normal stress amplitudes, and their products with the mean stresses, are averaged
over all plane orientations. The four weights are fixed by four fatigue strengths: fully reversed
tension and torsion, and repeated ($R=0$) tension and torsion. Proposed, from a data-driven study,
for the high-cycle fatigue limit of metals under multiaxial loading with mean stresses.

\block{Inputs}
\begin{center}\small
\begin{tabular}{@{}llp{100mm}@{}}
\toprule
Symbol & Unit & Meaning\\\midrule
$\n$, $S^2$ & -- & plane normal and the unit sphere of normals\\
$\tau_a$, $\tau_m$ & MPa & shear stress amplitude and mean on the plane\\
$\sigma_{n,a}$, $\sigma_{n,m}$ & MPa & normal stress amplitude and mean on the plane\\
$s_{-1}=\fm$, $t_{-1}=\tm$ & MPa & fully reversed axial and torsion fatigue strengths (amplitudes)\\
$s_0$, $t_0$ & MPa & repeated ($R=0$) axial and torsion strengths, maximum stress (amplitude and mean $s_0/2$)\\
$\sigma_u$ & MPa & ultimate tensile strength (estimate of $s_0$)\\
$\kappa=s_{-1}/t_{-1}$ & -- & limit ratio\\
$a,b,c,d$ & -- & weights, Eq.~\eqref{eq:bp-coef}\\
$\sigma_a$, $\sigma_m$ & MPa & uniaxial amplitude and mean of the worked example\\
$\sigma_{eq,BP}$ & MPa & equivalent stress, Eq.~\eqref{eq:bp-main}\\
$\mathrm{DP}$ & -- & fatigue index $\sigma_{eq,BP}/s_{-1}$\\
\bottomrule
\end{tabular}
\end{center}

\block{Formulation} The equivalent stress is the square root of the orientation average of a
quadratic form of the plane stresses (2026, Eq.~6):
\begin{equation}
\sigma_{eq,BP}^2=\frac{1}{4\pi}\int_{S^2}\left(a\,\tau_a^2+b\,\sigma_{n,a}^2+c\,\tau_a\tau_m+d\,\sigma_{n,a}\sigma_{n,m}\right)dS ,
\qquad \sigma_{eq,BP}\le s_{-1},
\label{eq:bp-main}
\end{equation}
\begin{equation}
a=\tfrac92\left(\kappa^2-\tfrac43\right),\quad
b=3\,(3-\kappa^2),\quad
c=\tfrac52\,\kappa^2\left[\Bigl(\frac{2t_{-1}}{t_0}\Bigr)^2-1\right],\quad
d=5\left[\Bigl(\frac{2s_{-1}}{s_0}\Bigr)^2-1-\frac{2c}{15}\right].
\label{eq:bp-coef}
\end{equation}
The amplitude weights are non-negative only for $4/3\le\kappa^2\le3$. When $s_0$, $t_0$ are not
measured (2026, Eqs.~1--2),
\begin{equation}
s_0=\frac{4\,s_{-1}\,\sigma_u}{s_{-1}+2\sigma_u},\qquad \frac{4t_{-1}}{t_0}-\frac{2s_{-1}}{s_0}=1 .
\label{eq:bp-r0}
\end{equation}
For uniaxial loading $\langle\cos^4\theta\rangle=1/5$ and $\langle\cos^2\theta\sin^2\theta\rangle=2/15$,
so
\begin{equation}
\sigma_{eq,BP}^2=\sigma_a^2+\sigma_a\sigma_m\left[\Bigl(\frac{2s_{-1}}{s_0}\Bigr)^2-1\right],
\label{eq:bp-uni}
\end{equation}
which equals $s_{-1}^2$ at $s_{-1}$ (fully reversed) and at $s_0$ ($\sigma_a=\sigma_m=s_0/2$).

\block{Constants and calibration} $a,b,c,d$ in closed form from $s_{-1}$, $t_{-1}$, $s_0$, $t_0$, the
$R=0$ strengths measured or estimated by Eq.~\eqref{eq:bp-r0}. The repository takes them from the
workbook \texttt{Sigm1}, \texttt{Taum1}, \texttt{Sig0}, \texttt{Tau0} (and \texttt{Su} for the
estimate); nothing is fitted except the master-curve constants on the uniaxial $R=-1$ design tests.

\block{Life} None: a fatigue-limit criterion. The repository reads finite lives on the master curve
of the index $\sigma_{eq,BP}/s_{-1}$ (Section~\ref{sec:mastercurve}).

\block{Repository choices}
\begin{itemize}
\item $\tau_a$ is the radius of the minimum circumscribed circle of the shear path (the major
  semi-axis of the single-harmonic ellipse); $\tau_m$ is the magnitude of the mean shear vector.
\item 24$\times$48 Gauss--Legendre/azimuth quadrature instead of the MATLAB angular grid: agreement
  within discretisation error, not bitwise.
\item A negative integral is clipped to 0. The admissible range $4/3\le\kappa^2\le3$ is not enforced.
\item Missing \texttt{Sig0} and/or \texttt{Tau0} are estimated by Eq.~\eqref{eq:bp-r0}; this needs
  \texttt{Su}.
\item Single-harmonic loading only; no \texttt{\_ext} variant (no plane is selected).
\end{itemize}

\block{Worked example} Synthetic material: $s_{-1}=300$, $t_{-1}=220$, $\sigma_u=800\MPa$
($\kappa^2=1.86$); Eq.~\eqref{eq:bp-r0} gives $s_0=505.3\MPa$. The index is 1 in reversed tension at
$s_{-1}$, in reversed torsion at $t_{-1}$ and at the $R=0$ limit. Tension $\sigma_a=200$ with
$\sigma_m=150\MPa$: index 0.7623, equal to Eq.~\eqref{eq:bp-uni}. Out of phase, $\sigma_{xx}=200\sin\omega t$
with $\tau_{xy}=100\cos\omega t$: index 0.7859, unchanged by a shift of the time origin.

\block{Verification} Checked: the four loading limits; phase-origin invariance; the $R=0$ estimates;
the closed forms of the 2026 Eq.~(17) for equibiaxial and equitriaxial amplitudes. Table~4 of the
2026 paper (statistics of the 18 Table~2 cases, static normal mean with alternating shear) is
reproduced from the 17 computable cases (case SE gives neither $s_0$ nor $t_0$ nor $\sigma_u$): Min
$-8.0\%$, Max $16.2\%$, SD $7.8\%$, Range $24.2\%$; BP and CAIM are identical there, as the paper
states. Its PM10 $=79.0\%$ is not a multiple of $1/18$ (misprint; $83.3\%$ follows from its other
values). The 2024 Eq.~(16) carries $1/(4\pi)$, the 2026 Eq.~(6) prints the integral without it;
only the sphere average reproduces the uniaxial limits and the closed forms of Eq.~(17), so the 2026
omission is a misprint. Not reproduced: the database results of the 2024 paper.
